% titel: Blatt 1, A5/A6: Schema der Materie / Teilchem-Tabelle
% skript: Mortimer Kap. 2 (Materie-Klassifikation)

% Teil I: Werkzeuge, kein Lösungsweg.

Schematische Einordnung der Materie und Zuordnungstabelle ausfüllen.

\begin{werkzeugkasten}[Materie-Schema]
Die Einteilung folgt: \textbf{Materie} $\rightarrow$ \emph{reine Stoffe} (Elemente,
Verbindungen) und \emph{Gemische} (homogen/heterogen). Jede Kategorie mit
typischen Vertretern:
\begin{itemize}
  \item Elemente: Atome einer Sorte (z. B.\ O$_2$, Fe).
  \item Verbindungen: chemisch gebundene Elemente (H$_2$O, CO$_2$).
  \item Homogene Gemische: Lösung, Gasgemisch, Legierung.
  \item Heterogene Gemische: Suspension, Emulsion, Aerosol.
\end{itemize}
\textit{Skript:} Mortimer Kap. 1.3 (Klassifikation der Materie).
\end{werkzeugkasten}

\begin{achtung}[Darauf achten]
\begin{itemize}
  \item Die Abgrenzung ``reiner Stoff vs.\ Gemisch" über die chemische
        Zusammensetzung, nicht über den Aggregatzustand.
  \item Verbindung $\neq$ Element: Verbindungen lassen sich chemisch zerlegen.
  \item Die konkreten Schemaboxen aus dem Blatt übernehmen (hier nicht
        ablesbar) und den Formalismus aus dem Skript eintragen.
\end{itemize}
\end{achtung}
